\section{Inversions, length and Lehmer codes}

Throughout this section, let $n \in \mathbb{N}$.
We write $S_n$ for the symmetric group on $[n] = \{1,2,\ldots,n\}$.

\subsection{Inversions and lengths}

\begin{definition}
\label{def.perm.invs}
\lean{AlgebraicCombinatorics.Perm.inv, AlgebraicCombinatorics.Perm.invCount}
\leantarget
\leanok
Let $n\in\mathbb{N}$ and $\sigma\in S_{n}$. \medskip

\textbf{(a)} An \emph{inversion} of $\sigma$ means a pair $\left(i,j\right)$
of elements of $\left[n\right]$ such that $i<j$ and $\sigma\left(i\right)
>\sigma\left(j\right)$. \medskip

\textbf{(b)} The \emph{length} (also known as the \emph{Coxeter length}) of
$\sigma$ is the \# of inversions of $\sigma$. It is called $\ell\left(
\sigma\right)$.
\end{definition}

\begin{lemma}
\label{lem.perm.inv_one}
\lean{AlgebraicCombinatorics.Perm.inv_one, AlgebraicCombinatorics.Perm.invCount_one}
\leanhelper
\leanok
The identity permutation $\operatorname{id} \in S_n$ has no inversions:
$\operatorname{inv}(\operatorname{id}) = \emptyset$ and $\ell(\operatorname{id}) = 0$.
\end{lemma}

\begin{proof}
\leanok
The identity preserves the natural order: if $i < j$ then $\operatorname{id}(i) = i < j = \operatorname{id}(j)$,
so no pair $(i,j)$ with $i < j$ satisfies $\operatorname{id}(i) > \operatorname{id}(j)$.
\end{proof}

\begin{proposition}
\label{prop.perm.lengths-k-small-k}
\lean{AlgebraicCombinatorics.Perm.invCount_le_choose, AlgebraicCombinatorics.Perm.invCount_eq_zero_iff, AlgebraicCombinatorics.Perm.card_invCount_eq_zero, AlgebraicCombinatorics.Perm.invCount_eq_choose_iff, AlgebraicCombinatorics.Perm.invCount_longestElement, AlgebraicCombinatorics.Perm.card_invCount_eq_choose, AlgebraicCombinatorics.Perm.card_invCount_eq_one, AlgebraicCombinatorics.Perm.card_invCount_eq_two, AlgebraicCombinatorics.Perm.card_invCount_symm}
\leantarget
\leanok
Let $n\in\mathbb{N}$. \medskip

\textbf{(a)} For any $\sigma\in S_{n}$, we have $\ell\left(\sigma\right)
\in\left\{0,1,\ldots,\binom{n}{2}\right\}$.
\lean{AlgebraicCombinatorics.Perm.invCount_le_choose}
\leanok
\medskip

\textbf{(b)} We have
\[
\left(\text{\# of }\sigma\in S_{n}\text{ with }\ell\left(\sigma\right)
=0\right) = 1.
\]
Indeed, the only permutation $\sigma\in S_{n}$ with $\ell\left(
\sigma\right)=0$ is the identity map $\operatorname{id}$.
\lean{AlgebraicCombinatorics.Perm.invCount_eq_zero_iff, AlgebraicCombinatorics.Perm.card_invCount_eq_zero}
\leanok
\medskip

\textbf{(c)} We have
\[
\left(\text{\# of }\sigma\in S_{n}\text{ with }\ell\left(\sigma\right)
=\binom{n}{2}\right) = 1.
\]
Indeed, the only permutation $\sigma\in S_{n}$ with $\ell\left(
\sigma\right)=\binom{n}{2}$ is the permutation with OLN $n\left(
n-1\right)\left(n-2\right)\cdots21$, often called~$w_{0}$.
\lean{AlgebraicCombinatorics.Perm.invCount_eq_choose_iff, AlgebraicCombinatorics.Perm.invCount_longestElement, AlgebraicCombinatorics.Perm.card_invCount_eq_choose}
\leanok
\medskip

\textbf{(d)} If $n\geq1$, then
\[
\left(\text{\# of }\sigma\in S_{n}\text{ with }\ell\left(\sigma\right)
=1\right) = n-1.
\]
Indeed, the only permutations $\sigma\in S_{n}$ with $\ell\left(
\sigma\right)=1$ are the simple transpositions $s_{i}$ with $i\in\left[
n-1\right]$.
\lean{AlgebraicCombinatorics.Perm.card_invCount_eq_one}
\leanok
\medskip

\textbf{(e)} If $n\geq2$, then
\[
\left(\text{\# of }\sigma\in S_{n}\text{ with }\ell\left(\sigma\right)
=2\right) = \dfrac{\left(n-2\right)\left(n+1\right)}{2}.
\]
Indeed, the only permutations $\sigma\in S_{n}$ with $\ell\left(
\sigma\right)=2$ are the products $s_{i}s_{j}$ with $1\leq i<j<n$ as well as
the products $s_{i}s_{i-1}$ with $i\in\left\{2,3,\ldots,n-1\right\}$. If
$n\geq2$, then there are $\dfrac{\left(n-2\right)\left(n+1\right)}{2}$
such products (and they are all distinct).
\lean{AlgebraicCombinatorics.Perm.card_invCount_eq_two}
\leanok
\medskip

\textbf{(f)} For any $k\in\mathbb{Z}$, we have
\[
\left(\text{\# of }\sigma\in S_{n}\text{ with }\ell\left(\sigma\right)
=k\right) = \left(\text{\# of }\sigma\in S_{n}\text{ with }\ell\left(
\sigma\right)=\binom{n}{2}-k\right).
\]
\lean{AlgebraicCombinatorics.Perm.card_invCount_symm}
\leanok
\end{proposition}

\begin{proof}
\leanok

\textbf{(a)}
The set of inversions of $\sigma$ is a subset of the set of all pairs $(i,j)$ with $i < j$, and the latter has cardinality $\binom{n}{2}$.

\textbf{(b)}
If $\sigma$ has no inversions, then $\sigma$ is strictly monotone on $[n]$; a strictly monotone permutation on a finite set must be the identity (since if any $\sigma(i) > i$, the sum $\sum_j \sigma(j)$ would strictly exceed $\sum_j j$, contradicting bijectivity).
Conversely, $\operatorname{id}$ clearly has no inversions.

\textbf{(c)}
The longest element $w_0$ reverses the ordering, so $w_0(i) = n+1-i$,
making every pair $(i,j)$ with $i<j$ an inversion. Thus $\ell(w_0) = \binom{n}{2}$.
Conversely, if $\ell(\sigma) = \binom{n}{2}$ then $\operatorname{inv}(\sigma)$ equals the set of all ordered pairs,
so $\sigma$ is strictly antitone.
A strictly antitone permutation on $[n]$ must send $i$ to $n+1-i$, i.e., $\sigma = w_0$.

\textbf{(d)}
Each simple transposition $s_i = \operatorname{swap}(i, i+1)$ has exactly one inversion.
Conversely, if $\sigma$ has exactly one inversion $\{(a,b)\}$, then $b = a + 1$
(otherwise an element between $a$ and $b$ would create an additional inversion),
and $\sigma = \operatorname{swap}(a, b) = s_a$.

\textbf{(e)}
The permutations with exactly 2 inversions are exactly:
products $s_i s_j$ with $1 \le i < j < n$ (non-adjacent transpositions, contributing $\binom{n-2}{2}$ permutations)
and products $s_i s_{i-1}$ with $i \in \{2, \ldots, n-1\}$ (adjacent 3-cycles, contributing $n-2$ permutations).
These give $\binom{n-2}{2} + (n-2) = \frac{(n-2)(n+1)}{2}$ in total.

\textbf{(f)}
The bijection $\sigma \mapsto w_0 \sigma$ satisfies $\ell(w_0\sigma) = \binom{n}{2} - \ell(\sigma)$,
since $w_0$ swaps inversions and non-inversions.
\end{proof}

\subsection{Helpers for Proposition~\ref{prop.perm.lengths-k-small-k}}

\begin{lemma}
\label{lem.perm.longestElement}
\lean{AlgebraicCombinatorics.Perm.longestElement, AlgebraicCombinatorics.Perm.longestElement_mul_self, AlgebraicCombinatorics.Perm.longestElement_inv, AlgebraicCombinatorics.Perm.longestElement_apply}
\leanhelper
\leanok
The longest element $w_0 \in S_n$ is the reversal permutation $w_0(i) = n+1-i$.
It is an involution: $w_0^2 = \operatorname{id}$ and $w_0^{-1} = w_0$.
\end{lemma}

\begin{proof}
\leanok
Direct computation.
\end{proof}

\begin{lemma}
\label{lem.perm.longestElement_eq_revPerm}
\lean{AlgebraicCombinatorics.Perm.longestElement_eq_revPerm}
\leanhelper
\leanok
The longest element $w_0 \in S_n$ equals the standard reversal
permutation from Mathlib. That is, $w_0(i) = n - 1 - i$ for all $i \in [n]$.
\end{lemma}

\begin{proof}
\leanok
Both send $i$ to $n - 1 - i$; the equality follows by extensionality.
\end{proof}

\begin{lemma}
\label{lem.perm.inv_longestElement}
\lean{AlgebraicCombinatorics.Perm.inv_longestElement}
\leanhelper
\leanok
The set of inversions of the longest element $w_0$ is the set of all pairs $(i,j)$
with $i < j$. That is, every ordered pair is an inversion of $w_0$.
\end{lemma}

\begin{proof}
\leanok
Since $w_0(i) = n+1-i$, if $i < j$ then $w_0(i) = n+1-i > n+1-j = w_0(j)$.
\end{proof}

\begin{definition}
\label{def.perm.nonInv}
\lean{AlgebraicCombinatorics.Perm.nonInv}
\leanhelper
\leanok
The \emph{non-inversions} of a permutation $\sigma \in S_n$ are the pairs $(i,j)$
with $i < j$ and $\sigma(i) < \sigma(j)$. These are exactly the ordered pairs
that are not inversions. The inversions and non-inversions partition the set of
all $\binom{n}{2}$ ordered pairs.
\end{definition}

\begin{lemma}
\label{lem.perm.invCount_inv}
\lean{AlgebraicCombinatorics.Perm.invCount_inv}
\leanhelper
\leanok
For any $\sigma \in S_n$, we have $\ell(\sigma^{-1}) = \ell(\sigma)$.
\end{lemma}

\begin{proof}
\leanok
The map $(a,b) \mapsto (\sigma^{-1}(b), \sigma^{-1}(a))$ is a bijection between
the inversions of $\sigma^{-1}$ and the inversions of $\sigma$.
\end{proof}

\begin{lemma}
\label{lem.perm.invCount_longestElement_mul}
\lean{AlgebraicCombinatorics.Perm.invCount_longestElement_mul', AlgebraicCombinatorics.Perm.invCount_longestElement_mul}
\leanhelper
\leanok
For any $\sigma \in S_n$, the inversions of $w_0 \sigma$ are exactly the non-inversions of $\sigma$, and
\[
\ell(w_0 \sigma) = \binom{n}{2} - \ell(\sigma).
\]
\end{lemma}

\begin{proof}
\leanok
A pair $(i,j)$ with $i < j$ is an inversion of $w_0 \sigma$ iff
$(w_0 \sigma)(i) > (w_0 \sigma)(j)$ iff $\sigma(i) < \sigma(j)$
(since $w_0$ reverses the order), which is the condition for $(i,j)$
to be a non-inversion of $\sigma$.
The cardinality identity follows since inversions and non-inversions
partition the $\binom{n}{2}$ ordered pairs.
\end{proof}

\begin{lemma}
\label{lem.perm.invCount_mul_le}
\lean{AlgebraicCombinatorics.Perm.invCount_mul_le}
\leanhelper
\leanok
For any $\sigma, \tau \in S_n$,
\[
\ell(\sigma \tau) \leq \ell(\sigma) + \ell(\tau).
\]
This is the triangle inequality for the inversion count.
\end{lemma}

\begin{proof}
\leanok
An inversion $(i,j)$ of $\sigma \tau$ (i.e., $i < j$ and $(\sigma\tau)(i) > (\sigma\tau)(j)$)
is either an inversion of $\tau$ (if $\tau(i) > \tau(j)$) or corresponds to an
inversion of $\sigma$ at $(\tau(i), \tau(j))$ (if $\tau(i) < \tau(j)$).
Thus $\operatorname{inv}(\sigma\tau) \subseteq \operatorname{inv}(\tau) \cup \tau^{-1}(\operatorname{inv}(\sigma))$,
and the result follows by cardinality.
\end{proof}

\subsection{Lehmer codes}

\begin{definition}
\label{def.perm.lehmer1}
\lean{AlgebraicCombinatorics.Perm.lehmerEntry, AlgebraicCombinatorics.Perm.Iic0, AlgebraicCombinatorics.Perm.Iic0Nat, AlgebraicCombinatorics.lehmerCodeSet, AlgebraicCombinatorics.card_lehmerCodeSet, AlgebraicCombinatorics.Perm.lehmerEntry_le, AlgebraicCombinatorics.Perm.lehmerCode, AlgebraicCombinatorics.Perm.lehmerCode_mem_lehmerCodeSet}
\leantarget
\leanok
Let $n\in\mathbb{N}$. \medskip

\textbf{(a)} For each $\sigma\in S_{n}$ and $i\in\left[n\right]$, we set
\begin{align*}
\ell_{i}\left(\sigma\right) := &
\left(\text{\# of all }j\in\left[n\right]
\text{ satisfying }i<j\text{ and }\sigma\left(i\right)
>\sigma\left(j\right)\right) \\
= & \left(\text{\# of all }j\in\left\{i+1,i+2,\ldots,n\right\}
\text{ such that }\sigma\left(i\right)>\sigma\left(j\right)\right).
\end{align*}
\lean{AlgebraicCombinatorics.Perm.lehmerEntry}
\leanok
\medskip

\textbf{(b)} For each $m\in\mathbb{Z}$, we let $\left[m\right]_{0}$ denote
the set $\left\{0,1,\ldots,m\right\}$. (This is an empty set when $m<0$.)
\lean{AlgebraicCombinatorics.Perm.Iic0, AlgebraicCombinatorics.Perm.Iic0Nat}
\leanok
\medskip

\textbf{(c)} We let $H_{n}$ denote the set
\begin{align*}
& \left[n-1\right]_{0}\times\left[n-2\right]_{0}\times\cdots
\times\left[n-n\right]_{0}\\
& = \left\{\left(j_{1},j_{2},\ldots,j_{n}\right)\in\mathbb{N}^{n}
\ \mid\ j_{i}\leq n-i\text{ for each }i\in\left[n\right]\right\}.
\end{align*}
This set $H_{n}$ has size $n!$.
\lean{AlgebraicCombinatorics.lehmerCodeSet, AlgebraicCombinatorics.card_lehmerCodeSet}
\leanok
\medskip

\textbf{(d)} Each $\sigma\in S_n$ and each $i\in[n]$ satisfy
$\ell_i(\sigma) \in \{0,1,\ldots,n-i\} = [n-i]_0$.
\lean{AlgebraicCombinatorics.Perm.lehmerEntry_le}
\leanok
\medskip

\textbf{(e)} We define the map
\begin{align*}
L:S_{n} &\rightarrow H_{n},\\
\sigma &\mapsto\left(\ell_{1}\left(\sigma\right),\ell_{2}\left(
\sigma\right),\ldots,\ell_{n}\left(\sigma\right)\right).
\end{align*}
If $\sigma\in S_{n}$ is a permutation, then the $n$-tuple $L\left(\sigma\right)$
is called the \emph{Lehmer code} (or just the \emph{code}) of $\sigma$.
\lean{AlgebraicCombinatorics.Perm.lehmerCode, AlgebraicCombinatorics.Perm.lehmerCode_mem_lehmerCodeSet}
\leanok
\end{definition}

\begin{proposition}
\label{prop.perm.lehmer.l}
\lean{AlgebraicCombinatorics.Perm.invCount_eq_sum_lehmerCode}
\leantarget
\leanok
Let $n\in\mathbb{N}$ and $\sigma\in S_{n}$. Then,
$\ell\left(\sigma\right)=\ell_{1}\left(\sigma\right)+\ell_{2}\left(
\sigma\right)+\cdots+\ell_{n}\left(\sigma\right)$.
\end{proposition}

\begin{proof}
\leanok
The inversions of $\sigma$ are pairs $(i,j)$ with $i<j$ and $\sigma(i)>\sigma(j)$.
The Lehmer entry $\ell_i(\sigma)$ counts the $j>i$ with $\sigma(i)>\sigma(j)$.
Hence the sum $\sum_i \ell_i(\sigma)$ partitions the set of inversions
by first coordinate, so equals $\ell(\sigma)$.
\end{proof}

\begin{theorem}
\label{thm.perm.lehmer.bij}
\lean{AlgebraicCombinatorics.Perm.lehmerCode_bijective, AlgebraicCombinatorics.Perm.lehmerCode_injective}
\leantarget
\leanok
Let $n\in\mathbb{N}$. Then, the map $L:S_{n}\rightarrow H_{n}$ is a bijection.
\end{theorem}

\begin{proof}[First proof of Theorem \ref{thm.perm.lehmer.bij} (sketched).]
\proves{thm.perm.lehmer.bij}
\leanok
Let $\sigma\in S_{n}$, $i\in[n]$.
We can rewrite
\begin{align}
\ell_{i}\left(\sigma\right)
&= \left(\text{\# of elements of }
\left[n\right]\setminus\left\{\sigma\left(1\right),\ldots,\sigma\left(i-1\right)\right\}
\text{ that are smaller than }\sigma\left(i\right)\right).
\label{eq.thm.perm.lehmer.bij.1st.1}
\end{align}
Therefore,
\begin{align}
\sigma\left(i\right) &= \left(\text{the }\left(\ell_{i}\left(\sigma\right)+1\right)\text{-st
smallest element of }
\left[n\right]\setminus\left\{\sigma\left(1\right),\ldots,\sigma\left(i-1\right)\right\}\right).
\label{eq.thm.perm.lehmer.bij.1st.2}
\end{align}
This allows recovering $\sigma(i)$ from $\ell_i(\sigma)$ and $\sigma(1),\ldots,\sigma(i-1)$
by induction on $i$, proving that $L$ is injective.

Conversely, given any $\mathbf{j} = (j_1,\ldots,j_n) \in H_n$,
define $M(\mathbf{j})$ to be the map $\sigma:[n]\to[n]$ whose values are determined by
\[
\sigma(i) = (\text{the }(j_i+1)\text{-st smallest element of }
[n]\setminus\{\sigma(1),\ldots,\sigma(i-1)\}).
\]
This is well-defined (since $j_i \le n - i$) and gives a permutation
(since $\sigma(i) \notin \{\sigma(1),\ldots,\sigma(i-1)\}$).
The maps $L$ and $M$ are mutually inverse, so $L$ is bijective.
\end{proof}

\begin{lemma}
\label{lem.perm.lehmerEntry_eq_card_filter}
\lean{AlgebraicCombinatorics.Perm.lehmerEntry_eq_card_filter}
\leanhelper
\leanok
For $\sigma \in S_n$ and $i \in [n]$,
\[
\ell_i(\sigma) = \left|\left\{ v < \sigma(i) : v \notin \sigma(\{1,\ldots,i-1\}) \right\}\right|.
\]
\end{lemma}

\begin{proof}
\leanok
The elements $j > i$ with $\sigma(j) < \sigma(i)$ are in bijection
(via $j \mapsto \sigma(j)$) with elements $v < \sigma(i)$ not in the image
$\sigma(\{1,\ldots,i-1\})$, since these are precisely
$[n] \setminus \{\sigma(1),\ldots,\sigma(i)\}$ intersected with $\{v : v < \sigma(i)\}$.
\end{proof}

\begin{lemma}
\label{lem.perm.lehmerEntry_lt}
\lean{AlgebraicCombinatorics.Perm.lehmerEntry_lt}
\leanhelper
\leanok
For $\sigma \in S_n$ and $i \in [n]$,
\[
\ell_i(\sigma) < n - i.
\]
This strict bound (compared to the weak bound $\ell_i(\sigma) \leq n - i$
from Definition~\ref{def.perm.lehmer1} part~(d)) holds because
$\ell_i(\sigma)$ counts elements among $\{i+1, \ldots, n-1\}$ where
$\sigma(j) < \sigma(i)$, and there are only $n - 1 - i$ such elements.
\end{lemma}

\begin{proof}
\leanok
The set $\{j : j > i \text{ and } \sigma(j) < \sigma(i)\}$ is a subset
of $\{j : j > i\}$, which has cardinality $n - 1 - i < n - i$.
\end{proof}

\subsection{Lexicographic order}

\begin{definition}
\label{def.perm.lehmer.lex-ord}
\lean{AlgebraicCombinatorics.lexLt}
\leantarget
\leanok
Let $\left(a_{1},a_{2},\ldots,a_{n}\right)$
and $\left(b_{1},b_{2},\ldots,b_{n}\right)$ be two $n$-tuples of integers.
We say that
\[
\left(a_{1},a_{2},\ldots,a_{n}\right) <_{\operatorname{lex}}\left(
b_{1},b_{2},\ldots,b_{n}\right)
\]
if and only if
\begin{itemize}
\item there exists some $k\in\left[n\right]$ such that $a_{k}\neq b_{k}$, and
\item the \textbf{smallest} such $k$ satisfies $a_{k}<b_{k}$.
\end{itemize}
The relation $<_{\operatorname{lex}}$ is a strict total order on $\mathbb{Z}^n$
(the lexicographic order).
\lean{AlgebraicCombinatorics.lexLt_irrefl, AlgebraicCombinatorics.lexLt_trans, AlgebraicCombinatorics.lexLt_asymm, AlgebraicCombinatorics.lexLt_strictTotalOrder}
\leanok
\end{definition}

\begin{proposition}
\label{prop.perm.lehmer.lex-ord.total}
\lean{AlgebraicCombinatorics.lexLt_trichotomous}
\leantarget
\leanok
If $\mathbf{a}$ and $\mathbf{b}$ are two
distinct $n$-tuples of integers, then we have either $\mathbf{a}
<_{\operatorname{lex}}\mathbf{b}$ or $\mathbf{b}<_{\operatorname{lex}}
\mathbf{a}$.
\end{proposition}

\begin{proof}
\leanok
Since $\mathbf{a} \neq \mathbf{b}$, there exists some index where they differ.
By well-founded induction on $\operatorname{Fin} n$, let $k$ be the minimal such index.
Then $a_i = b_i$ for all $i < k$, and $a_k \neq b_k$.
If $a_k < b_k$ then $\mathbf{a} <_{\operatorname{lex}} \mathbf{b}$;
otherwise $b_k < a_k$ and $\mathbf{b} <_{\operatorname{lex}} \mathbf{a}$.
\end{proof}

\begin{proposition}
\label{prop.perm.lehmer.lex}
\lean{AlgebraicCombinatorics.Perm.lehmerCode_preserves_lexLt}
\leantarget
\leanok
Let $\sigma\in S_{n}$ and $\tau\in S_{n}$ be such that
\[
\left(\sigma\left(1\right),\sigma\left(2\right),\ldots,\sigma\left(
n\right)\right) <_{\operatorname{lex}}\left(\tau\left(1\right)
,\tau\left(2\right),\ldots,\tau\left(n\right)\right).
\]
Then,
\[
\left(\ell_{1}\left(\sigma\right),\ell_{2}\left(\sigma\right)
,\ldots,\ell_{n}\left(\sigma\right)\right) <_{\operatorname{lex}}\left(
\ell_{1}\left(\tau\right),\ell_{2}\left(\tau\right),\ldots,\ell
_{n}\left(\tau\right)\right).
\]
(In other words, $L\left(\sigma\right) <_{\operatorname{lex}} L\left(
\tau\right)$.)
\end{proposition}

\begin{proof}[Proof of Proposition \ref{prop.perm.lehmer.lex} (sketched).]
\proves{prop.perm.lehmer.lex}
\leanok
The assumption says there exists a smallest $k \in [n]$ with
$\sigma(k) \neq \tau(k)$ and $\sigma(k) < \tau(k)$.
For $i < k$: $\sigma(i) = \tau(i)$, so the images
$\sigma(\{1,\ldots,i-1\})$ and $\tau(\{1,\ldots,i-1\})$ agree.
Since $\sigma(i) = \tau(i)$ as well, we get $\ell_i(\sigma) = \ell_i(\tau)$.

Let $Z = [n] \setminus \{\sigma(1),\ldots,\sigma(k-1)\}
= [n] \setminus \{\tau(1),\ldots,\tau(k-1)\}$.
Then $\ell_k(\sigma)$ counts elements of $Z$ smaller than $\sigma(k)$,
and $\ell_k(\tau)$ counts elements of $Z$ smaller than $\tau(k)$.
Since $\sigma(k) < \tau(k)$ and $\sigma(k) \in Z$,
every element of $Z$ smaller than $\sigma(k)$ is also smaller than $\tau(k)$,
and $\sigma(k)$ itself is smaller than $\tau(k)$ but not smaller than $\sigma(k)$.
Thus $\ell_k(\sigma) < \ell_k(\tau)$.

Combining, we get $L(\sigma) <_{\operatorname{lex}} L(\tau)$.
\end{proof}

\begin{proof}[Second proof of Theorem \ref{thm.perm.lehmer.bij} (sketched).]
\proves{thm.perm.lehmer.bij}
\leanok
We first show that $L$ is injective.
Let $\sigma$ and $\tau$ be two distinct permutations in $S_n$.
Their OLNs are distinct, so by Proposition \ref{prop.perm.lehmer.lex-ord.total}
we have either OLN$(\sigma) <_{\operatorname{lex}}$ OLN$(\tau)$ or vice versa.
In either case, Proposition \ref{prop.perm.lehmer.lex} gives
$L(\sigma) <_{\operatorname{lex}} L(\tau)$ or
$L(\tau) <_{\operatorname{lex}} L(\sigma)$.
In particular $L(\sigma) \neq L(\tau)$, so $L$ is injective.
Since $|S_n| = n! = |H_n|$ and $L$ is injective between finite sets of equal size,
the Pigeonhole Principle shows $L$ is bijective.
\end{proof}

\subsection{Generating function for lengths}

\begin{proposition}
\label{prop.perm.length.gf}
\lean{AlgebraicCombinatorics.length_generating_function}
\leantarget
\leanok
Let $n\in\mathbb{N}$. Then,
\begin{align*}
\sum_{\sigma\in S_{n}} x^{\ell\left(\sigma\right)}
&= \prod_{i=1}^{n-1}\left(1+x+x^{2}+\cdots+x^{i}\right) \\
&= \left(1+x\right)\left(1+x+x^{2}\right)\left(1+x+x^{2}+x^{3}\right)
\cdots\left(1+x+x^{2}+\cdots+x^{n-1}\right) \\
&= \left[n\right]_{x}!.
\end{align*}
\end{proposition}

\begin{proof}[Proof of Proposition \ref{prop.perm.length.gf}.]
\proves{prop.perm.length.gf}
\leanok
Each $\sigma\in S_{n}$ satisfies
\[
\ell\left(\sigma\right) = \ell_{1}\left(\sigma\right)+\ell_{2}\left(
\sigma\right)+\cdots+\ell_{n}\left(\sigma\right)
\]
by Proposition \ref{prop.perm.lehmer.l},
whence $x^{\ell(\sigma)} = \prod_{i=1}^{n} x^{\ell_i(\sigma)}$.
Therefore
\begin{align*}
\sum_{\sigma\in S_{n}} x^{\ell\left(\sigma\right)}
&= \sum_{\sigma\in S_{n}} \prod_{i=1}^{n} x^{\ell_i(\sigma)}
= \sum_{\left(j_{1},\ldots,j_{n}\right)\in H_{n}}
x^{j_{1}} x^{j_{2}} \cdots x^{j_{n}}
\end{align*}
where we substituted $(j_1,\ldots,j_n)$ for $L(\sigma)$, which is valid since
$L: S_n \to H_n$ is a bijection (Theorem \ref{thm.perm.lehmer.bij}).
Since $H_{n} = [n-1]_0 \times [n-2]_0 \times \cdots \times [0]_0$,
the product rule gives
\begin{align*}
\sum_{(j_1,\ldots,j_n) \in H_n} x^{j_1}\cdots x^{j_n}
&= \left(\sum_{j_1=0}^{n-1} x^{j_1}\right)
   \left(\sum_{j_2=0}^{n-2} x^{j_2}\right)
   \cdots
   \left(\sum_{j_n=0}^{0} x^{j_n}\right) \\
&= \prod_{i=1}^{n-1} \left(1+x+x^2+\cdots+x^i\right)
= [n]_x!.
\end{align*}
\end{proof}
